
\documentclass[aspectratio=169]{beamer}

\mode<presentation>
\usetheme{Boadilla}
\definecolor{NTHU1}{RGB}{127,16,132}
\definecolor{NTHU2}{RGB}{228,210,231}
\definecolor{blue}{RGB}{30,90,205}
\definecolor{red}{RGB}{213,94,0}
\definecolor{green}{RGB}{0,128,0}
\definecolor{charcoalgray}{RGB}{108,104,104}
\setbeamercolor{title}{fg=NTHU1}
\setbeamercolor{frametitle}{fg=NTHU1}
\setbeamercolor{block title}{bg=NTHU1, fg=white}
\setbeamercolor{block body}{bg=NTHU2}
\setbeamercolor{structure}{fg=NTHU1}
\setbeamercolor{item projected}{fg=white}
\setbeamercolor{item}{fg=NTHU1}
\setbeamercolor{subitem}{fg=NTHU1}
\setbeamercolor{section in toc}{fg=NTHU1}
\setbeamercolor{description item}{fg=NTHU1}
\setbeamercolor{caption name}{fg=NTHU1}
\setbeamercolor{button}{bg=NTHU1, fg=white}
\setbeamercolor{caption name}{fg=NTHU1}
\usepackage{graphics}
\usepackage{tikz}
\usepackage{amsmath}
\usepackage{bbm}
\usetikzlibrary{decorations.pathreplacing}
\usepackage{geometry}
\usepackage{booktabs}
\usepackage{bookmark}
\usepackage{multirow, makecell}
\usepackage{float}
\usepackage{fancyvrb}
\usepackage{caption}
\captionsetup[figure]{singlelinecheck = false, format= hang, justification=centering}
\captionsetup[subfigure]{singlelinecheck = false, format= hang, justification=raggedright}
\usepackage{subcaption}
\usepackage{adjustbox}
\usepackage{hyperref}
\usepackage{threeparttable}
\usepackage{appendixnumberbeamer} 
\usepackage[T1]{fontenc}
\usepackage[default]{lato} 
\usepackage{smartdiagram}
\smartdiagramset{
  back arrow disabled = true,
  module minimum width=1cm,
  module x sep = 3,
  uniform arrow color=true,
}
\usepackage{xcolor}
\usepackage{colortbl}
  \definecolor{light-gray}{gray}{0.99}

\newenvironment{wideitemize}{\itemize\addtolength{\itemsep}{10pt}}{\enditemize}
\newenvironment{wideenumerate}{\enumerate\addtolength{\itemsep}{10pt}}{\endenumerate}
\newenvironment{widedescription}{\description\addtolength{\itemsep}{10pt}}{\enddescription}
\hypersetup{
colorlinks=true,
linkcolor=NTHU1,
filecolor=green, 
urlcolor=blue,
}
\beamertemplatenavigationsymbolsempty
\setbeamercolor{author in head/foot}{bg=white, fg=NTHU1}
\setbeamercolor{title in head/foot}{bg=white, fg=NTHU1}
\setbeamercolor{date in head/foot}{bg=white, fg=NTHU1}
\setbeamercolor{section in head/foot}{bg=white, fg=NTHU1}
\setbeamercolor{headline}{bg=white}
\setbeamertemplate{footline}{
    \leavevmode%
    \hbox{%
        \begin{beamercolorbox}[wd=.333333\paperwidth,ht=2.25ex,dp=1ex,center]{date in head/foot}%
            \usebeamerfont{date in head/foot}%\insertdate
        \end{beamercolorbox}%
        \begin{beamercolorbox}[wd=.333333\paperwidth,ht=2.25ex,dp=1ex,center]{title in head/foot}%
            \usebeamerfont{title in head/foot}\insertshorttitle
        \end{beamercolorbox}%
        \begin{beamercolorbox}[wd=.333333\paperwidth,ht=2.25ex,dp=1ex,right]{date in head/foot}%
            \usebeamerfont{date in head/foot}\insertframenumber{}\hspace*{2em}
        \end{beamercolorbox}}%
        \vskip0pt%
    }
    %% May need to script the introduction!!!!!!
%\setbeamercolor{page number in head/foot}{fg=NTHU1}
\setbeamertemplate{section in toc}[sections numbered]
\setbeamertemplate{subsection in toc}{\leavevmode\leftskip=3em\rlap{\hskip-1.75em\inserttocsectionnumber.\inserttocsubsectionnumber}
\inserttocsubsection\par}
\setbeamerfont{subsection in toc}{size=\footnotesize}
%\setbeamertemplate{headline}{%
  %\begin{beamercolorbox}[ht=5.5ex]{section in head/foot}
    %\vskip2pt\insertnavigation{0.33\paperwidth}\vskip2pt
  %\end{beamercolorbox}%
%}
\newenvironment{transitionframe}{\setbeamercolor{background canvas}{bg=white}\setbeamertemplate{footline}{} \begin{frame}[noframenumbering]}{\end{frame}}

\makeatletter
\let\@@magyar@captionfix\relax
\makeatother

\title[]{Conflict and Peace} % Change this regularly
\subtitle[]{Week 12: State capacity -- Theoretical framework and weak states}
\author[Lee]{Seung-hun Lee }
\institute[]{Taipei School of Economics\\ National Tsing Hua University}

\date[]{November 21st, 2025}

\begin{document}
\begin{transitionframe}
\titlepage
\end{transitionframe}

%%% Color slides for section headers: Use for colloquium version (The ones bewteen \iffals and \fi)

\section{Overview of the topic}
\begin{frame}
\frametitle{Drastic differences in capabilities of states across different countries} 
Some of the poorest countries in the world are also  `weak states'
\begin{itemize}\setlength\itemsep{3pt}
  \item State capacity: Encompasses ability to support markets, provide public goods, and raise taxes from broad bases 
   \begin{itemize}\setlength\itemsep{3pt}
    \item Support markets: Enforce contracts, protect property rights
    \item Provide public goods: Education, health, but also peace (by monopolizing violence)
    \item Broad base taxes (income tax) are more rule-based, less burdensome for individuals, but require large investments for implementation 
    \item Narrow base taxes (royalties and tariffs) are easy to implement but could be corruptible and extracts particular group excessively 
  \end{itemize}
\item Weak states: States that cannot support basic economic function, raise revenues, deliver essential public services, and provide law and order
\end{itemize} \medskip 
State capacity is widely dispersed, but also clustered
\begin{itemize}\setlength\itemsep{3pt}
\item Weak states are also poorer and mired in internal conflicts
\item Developed countries: High income, functioning institutions, less internal conflicts
\end{itemize}
\end{frame}

\begin{frame}
\frametitle{Huge variation in state capacity across the globe}
\begin{figure}
  \includegraphics[keepaspectratio,width=0.7\textwidth]{fig1.png}
\end{figure}
\end{frame}

\begin{frame}
\frametitle{State capacity clustered with income and functioning of institutions}
\begin{figure}
  \begin{subfigure}{0.49\textwidth}
  \includegraphics[keepaspectratio,width=\textwidth]{fig2_1.png}
\caption{Income and capacities are clustered}
  \end{subfigure}
  \begin{subfigure}{0.49\textwidth}
  \includegraphics[keepaspectratio,width=\textwidth]{fig2_2.png}
  \caption{Capacities and fragilities index}
  \end{subfigure}
\end{figure}
\end{frame}

\begin{frame}
\frametitle{External wars increase capacity; internal wars hamper it}
\begin{figure}
  \begin{subfigure}{0.49\textwidth}
  \includegraphics[keepaspectratio,width=\textwidth]{fig3_1.png}
\caption{External wars positively correlate}  
\end{subfigure}
  \begin{subfigure}{0.49\textwidth}
  \includegraphics[keepaspectratio,width=\textwidth]{fig3_2.png}
  \caption{Internal wars negatively correlate}
  \end{subfigure}
\end{figure}
\end{frame}

\begin{frame}
\frametitle{Why care about state capacity, and open questions?}
Why do state capacities matter?
\begin{itemize}\setlength\itemsep{3pt}
\item Capable states can distribute fruits of economic development and prevent violence that stifles the progress
\begin{itemize}
\item Ability to tax, enforce law, and distribute services contributes to growth and peace
\end{itemize}
\item Ultimately relates to economic and institutional development: Many weak states are poor, authoritarian, and in internal conflict
\end{itemize}\medskip

We discuss open questions on state capacity with theories and empirical findings
\begin{itemize}\setlength\itemsep{3pt}
\item What forces explain the formation of capable states?
\item How do weak states perpetuate violence, corruption, and weak growth?
\item What explains the clustering of state institutions, violence, and income?
\end{itemize}
\end{frame}


\begin{transitionframe}
\frametitle{Roadmap for today}
\tableofcontents
\end{transitionframe}

\section{Conceptualizing state capacity}
\subsection{Besley, Persson (2009): The Origins of State Capacity: Property Rights, Taxation and Politics}
\begin{transitionframe}
\begin{center}
  \begin{huge}
    \textcolor{NTHU1}{%
Besley, Persson (2009):\\ The Origins of State Capacity: Property Rights, Taxation and Politics
    }
  \end{huge}
\end{center}
\end{transitionframe}

\begin{frame}
\frametitle{Q: What determines the state's investment in fiscal and legal capacity}
How are institutional capacity support markets and collect taxes developed?
\begin{itemize}\setlength\itemsep{3pt}
\item Historians note the development of capacity in response to external wars (Tilly 1990)
\item Social sciences (economics, especially) typically assume such capacities exist by default
\item In reality, institutional capacities differ starkly
\end{itemize}\medskip
This paper provides theoretical framework of state institutions as endogenous investments
\begin{itemize}\setlength\itemsep{3pt}
\item Individuals come from different groups (ruling vs nonruling, different wealth)
\item State collects taxes, protects property rights, and provide public goods
\item What determines taxation, public goods provision, and support for market activities
\item Why are there differences in investments in these capacities across different settings?
\begin{itemize}\setlength\itemsep{3pt}
\item Focus on high capacities of wealthy, democratic, and parliamentary democracies 
\end{itemize}
\item Provides grounds for studying how institutions evolve during development
\end{itemize}
\end{frame}

\begin{frame}
\frametitle{Setting up the framework: Definitions and layout} 
What makes up state capacity
\begin{itemize}\setlength\itemsep{3pt}
\item Legal capacity: Enforce contracts, regulate markets, and protect property rights
\item Fiscal capacity: Collect taxes to finance public goods and `redistribution'
\end{itemize}\medskip 
Setting the model up
\begin{itemize}
\item There are two time periods $s\in\{1,2\}$ and two groups of individuals  $J\in\{A,B\}$
\item Individuals have linear preferences over public goods and post-tax income
\item State (led by $A$ or $B$) protects property rights, collect taxes in both periods
\item State can also increase its legal and fiscal capacities at $s=1$
\item State maximizes (politically) weighted sum of individual utilities under budget and capacity constraints
\end{itemize}
\end{frame}

\begin{frame}
\frametitle{Setting up the framework: Individual utilities} 
Individual utilities (For given group $J\in\{A,B\}$)
\begin{itemize}\setlength\itemsep{3pt}
\item Gain utility from public good $G_s$ (by factor of $\alpha_s$) and income $Y_s^J$ (taxed by $t_s^J$)
\[
v_s^J=\alpha_s G_s + (1-t_s^J)Y_s^J(p_s^J) 
\]\vspace{-20pt}
\item Income gain is higher with increase in $p_s^J$, enforcement of property rights, $\left(\frac{\partial Y_s^J}{\partial p_s^J} >0\right)$
\item Taxes deduct from income but used to finance public goods
\end{itemize}\medskip 
Differences across groups
\begin{itemize}\setlength\itemsep{3pt}
\item Demographic share: $\beta^J$ of population in group $J\in\{A,B\}$ ($\beta^A+ \beta^B=1$)
\item Different protection, taxation (negative taxes possible), and income 
\item Public good is distributed commonly (no differences in level of public goods)
\item Assume that group $A$ is ruling at $s=1$, with chance of continued ruling in $s=2$ is $\gamma_A$
\end{itemize}\medskip 

\end{frame}

\begin{frame}
\frametitle{Setting up the framework: State objectives and constraints} 
Objective of the state: Politically weighted sum of individual utilities
\begin{itemize}
\item $\bar{\rho}$ and $\underline{\rho}$: Relative political weights on ruling and nonruling groups ($\bar{\rho}\beta_s^A+\underline{\rho}\beta_s^B=1$)
\[
V_s = \alpha_s G_s + \bar{\rho}\beta_s^A(1-t_s^A)Y_s^A(p_s^A)+\underline{\rho}\beta_s^B(1-t_s^B)Y_s^B(p_s^B)
\]\vspace{-20pt}
\begin{itemize}
\item $\bar{\rho}$ captures how $A$ values its own group relative to demographic share (1 or greater)
\item $\underline{\rho}$ captures how $A$ values nonruling groups (less or equal to 1)
\end{itemize}
\item At $s=1$, invest in legal ($\pi_s$) and fiscal capacity ($\tau_s$) at cost of $L(\pi_2-\pi_1)$ and $F(\tau_2-\tau_1)$
\end{itemize} \medskip
Constraints of the state
\begin{itemize}\setlength\itemsep{3pt}
\item Capacity constraints: $p_s^J \leq \pi_s$, $t_s^J \leq \tau_s$
\item Budget Constraints
\begin{itemize}\setlength\itemsep{3pt}
\item $s=1$: Capacity investment included $\to \sum_J \beta^J t_1^J Y_1^J=G_1 + L(\pi_2-\pi_1)+F(\tau_2-\tau_1)$ 
\item $s=2$: No new investment$\to \sum_J \beta^J t_2^J Y_2^J=G_2 $ 
\end{itemize}
\end{itemize}
\end{frame}

\begin{frame}
\frametitle{Policy choices in utilitarian vs politically-motivated settings} 
Policy choices on $p_s^J$ and $t_s^J$ in utiltarian governments ($\bar{\rho}=\underline{\rho}=1$)
\begin{itemize}\setlength\itemsep{3pt}
\item $p_s^J$ increases indirect utilities: Provide to the fullest possible $p_s^J=\pi_s$ 
\item $t_s^J$ works in both ways - incomes are taxed, but used for public goods
\item $\alpha_s\geq1$: Tax to the fullest possible ($t_s^J=\tau_s$), all revenues to $G_s$ and investments 
\item $\alpha_s<1$: $G_s=0$, smallest possible taxes (0 in $s=2$, exact cost of investments in $s=1$)
\end{itemize}\medskip
Politically-motivated objective maximization ($\bar{\rho}>1>\underline{\rho}$)
\begin{itemize}\setlength\itemsep{3pt}
\item No different logic in legal capacity: Higest protection possible
\item Public goods are provided only if ruling class values it sufficiently ($\alpha_s\geq\bar{\rho}$)
\item If $\alpha<\bar{\rho}$:  $G_s=0$, and taxes could be `redistributed' from nonruling to ruling class
\begin{itemize}
  \item Ruling group gains $\bar{\rho}$ from `redistribution', nonruling loses $\underline{\rho}$ from this
  \item With $\bar{\rho}-\underline{\rho}>0$ beneficial to extract (from nonruling group standpoint)
\item Implies that in this setting, people have different interests towards fiscal capacity
\end{itemize}
\end{itemize}
\end{frame}

\begin{frame}
\frametitle{Investing in state capacity} 
Investments in fiscal ($\tau$) and legal ($\pi$) capacities maximizes expected social utility
\begin{itemize}\setlength\itemsep{3pt}
\item Factors: Wealth ($Y_s^J$), Value of public goods ($\alpha_s$), Regime stability ($\gamma_A$), representativeness ($\bar{\rho}-\underline{\rho}$)
\item Legal capacity is desirable in any case
\item When $\alpha_s<\bar{\rho}$, high $\tau$ undesirable for non-ruling group $\to$ Underinvest in capacities
\begin{itemize}
  \item In other cases, no conflict over the desirability of fiscal capacity
\end{itemize}
\end{itemize}\medskip
Key findings (in theory and empirics)
\begin{itemize}\setlength\itemsep{3pt}
\item In most cases, legal and fiscal capacities complement each other
\begin{itemize}
  \item $\pi\Uparrow\to$ More revenues for public goods $\to\tau\Uparrow\to$ More funds to increase protection$\to\pi\Uparrow$
\end{itemize}
\item What factors increase investments in state capacities?
\begin{itemize}
  \item High $Y_s^J$, High $\alpha_s$, High $\gamma_A$, Small $\bar{\rho}-\underline{\rho}$
  \item Mismatch in political power and economic influence $\to$ Fiscal capacity underinvested
\end{itemize}
\end{itemize}
\end{frame}

\begin{frame}
\frametitle{Theoretical prediction matches well with observed correlates} 
\begin{figure}
  \includegraphics[keepaspectratio, width=0.7\textwidth]{bp2009_t1.png}
\end{figure}
\end{frame}

\begin{frame}
\frametitle{Theoretical prediction matches well with observed correlates} 
\begin{figure}
  \includegraphics[keepaspectratio, width=0.7\textwidth]{bp2009_t2.png}
\end{figure}
\end{frame}

\subsection{Gennaioli, Voth (2015): State Capacity and Military Conflict}
\begin{transitionframe}
\begin{center}
  \begin{huge}
    \textcolor{NTHU1}{%
Gennaioli, Voth (2015):\\ State Capacity and Military Conflict
    }
  \end{huge}
\end{center}
\end{transitionframe}

\begin{frame}
\frametitle{Q: What explains the rise of centralized states in Europe?} 
Centralized states are relatively recent concept
\begin{itemize}\setlength\itemsep{3pt}
\item Before: No professional bureaucracy, limited tax powers, armies of mercenaries
\item Europe began to form centralized fiscal appartus from 1500s
\item Key driver: External wars - but really?
\begin{itemize}\setlength\itemsep{3pt}
\item Many wars pre-1500s, and many states (Spain, Austria) late in developing fiscal capacity
\item War may not be for common interests (personal glory)
\item Military capacity $\to$ war (reverse causality feasible)
\end{itemize}
\end{itemize}\medskip
This paper develops and tests the model of different types of wars in state-building
\begin{itemize}\setlength\itemsep{3pt}
\item War differs by how much money matters for military success
\item Cohesiveness of the pre-war state also important
\item Tests this on data of 374 battles in Europe since 1500
\item With rising importance of money in wars, cohesive states developed more
\end{itemize}
\end{frame}



\begin{frame}
\frametitle{Historigal background and theoretical framework} 
European wars and fiscal capacity
\begin{itemize}\setlength\itemsep{3pt}
\item Role of money became important around 1500s (but not before)
\begin{itemize}\setlength\itemsep{3pt}
\item Military Revolution: Standing army and new technology (gunpowder) made money crucial
\item Divergence of military powers followed (ENG, FRA vs SPA, AUT, POL)
\end{itemize}
\item Some states increased taxation since 1500s, but large cross-country variations
\begin{itemize}\setlength\itemsep{3pt}
  \item 10-fold increase for England, Decrease for Ottoman Empire and Poland
  \item Reliance on domain income $\to$ systemized tax collection (income, farming-based)
\end{itemize}\medskip
\end{itemize}\medskip
Theoretical model: Centralized vs decentralized state without external wars
\begin{itemize}\setlength\itemsep{3pt}
  \item There are local, market, and home activities (not taxed)
  \item Centralized: Local and market activities taxed by central ruler
  \item Decentralized: Market activities taxed across different localities (Thus overtaxed)
  \begin{itemize}\setlength\itemsep{3pt}
  \item Local production only taxed at base. Local powerholder collects taxes and gets to keep it
  \item Cost of centralization: Need to fend off local powerholders who lose tax revenues
\end{itemize}\medskip
\end{itemize}
\end{frame}

\begin{frame}
\frametitle{External wars and centralized state capacity} 
How does external war come into play?
\begin{itemize}\setlength\itemsep{3pt}
\item War absorbs revenues of warring rulers + Redistributes revenues from loser to winner
\item Winning depends on military strength (determined by manpower and expenses)
\item War threat affects centralization incentives in two ways
\begin{itemize}\setlength\itemsep{3pt}
\item Increase centralization: Rising revenues makes it more likely to win 
\item Decrease centralization: Revenues spent on war, and risk of larger losses after war 
\end{itemize}
\item Determinant of divergent patterns of state-building
\begin{itemize}\setlength\itemsep{3pt}
\item Higher money sensitivity: Increase incentive to centralize
\item Cohesiveness: External war becomes common threat $\to$ more incentive to centralize
\item Countries less likely to win (divided to begin with) unlikely to centralize 
\end{itemize}
\item When money does not matter: Less incentive to centralize 
 \end{itemize}
\end{frame}

\begin{frame}
\frametitle{Data and empirical strategy} 
Variables and data sources (Jacques (2007), Karaman, Pamuk (2010))
\begin{itemize}\setlength\itemsep{3pt}
\item Data on 374 European wars as external wars
\item Fiscal capacity: Tax revenues per capita divided by wage (mimic share of taxes)
\item Cohesion measured based on number of independent predecessor states, ethnic fractionalization measure, and surface area
\end{itemize}\medskip
Empirical strategy seeks to estimate correlates between money sensitivity, cohesion, and fiscal capacity
\begin{itemize}\setlength\itemsep{3pt}
\item Primarily use OLS (along with country fixed effects)

\end{itemize}\medskip

\end{frame}

\begin{frame}
\frametitle{Empirical finding: Richer countries won more wars after 1650} 
\begin{figure}
\includegraphics[keepaspectratio, width=0.58\textwidth]{gv2015_f1.png}
\end{figure}
\end{frame}

\begin{frame}
\frametitle{Sufficient cohesion and money sensitivity matters} 
\begin{figure}
\includegraphics[keepaspectratio, width=0.5\textwidth]{gv2015_t1.png}
\end{figure}
\end{frame}

\begin{frame}
\frametitle{Takeaway and discussion} 
Not all external wars affect state capacity equally
\begin{itemize}\setlength\itemsep{3pt}
\item Early history wars did not contribute to fiscal capacity
\item When revenues mattered in winning wars, it increased incentive to centralize
\item The effect is augmented with initial cohesion
\begin{itemize}
\item These patterns explain why centralization patterns diverged even in Europe
\end{itemize}
\end{itemize}\medskip
Remaining questions
\begin{itemize}\setlength\itemsep{3pt}
\item With decreased frequencies of war, what explains development of fiscal capacity?
\begin{itemize}\setlength\itemsep{3pt}
\item Ability to implement broad-based tax still key
\item Potential role for digitalized tax system (Pomeranz 2015; Okunogbe, Tourek 2024)
\end{itemize}
\item Maybe investment in military capacity is a result of possible external war
\begin{itemize}\setlength\itemsep{3pt}
\item If military investment = state capacity, then Besley, Persson (2009) causal logic (war $\to$ invest capacity) still applicable
\end{itemize}
\end{itemize}\medskip
\end{frame}

\subsection{Besley (2020): State Capacity, Reciprocity, and Social Contract}
\begin{transitionframe}
\begin{center}
  \begin{huge}
    \textcolor{NTHU1}{%
Besley (2020):\\ State Capacity, Reciprocity, and Social Contract
    }
  \end{huge}
\end{center}
\end{transitionframe}

\begin{frame}
\frametitle{Q: How does civic culture contribute to fiscal capacity?}
Increased tax intake enabled governments to spend on essential services 
\begin{itemize}\setlength\itemsep{3pt}
\item Not just on national defense, but also health care, education, income transfers etc
\item Switches in public spending requires economic, political, and social changes
\begin{itemize}\setlength\itemsep{3pt}
\item Should develop capacity to tax and raise revenues
\item But also, cultivate civic culture that increases tax \textit{compliance}
\end{itemize}
\end{itemize}\medskip
This paper extends the Besley, Persson (2009) model
\begin{itemize}\setlength\itemsep{3pt}
\item  By incorporating the role of civic culture, reciprocity, and social contract
\item Relative payoffs of civic-minded citizens and materialistic citizens matter
\item Investigates how civic culture can develop in the long-run
\item Ultimately answers how state capacity can persist with interactions with citizens
\end{itemize}
\end{frame}

\begin{frame}
\frametitle{Motivating findings: Different attitudes to tax compliance} 
\begin{figure}
\includegraphics[keepaspectratio, width=0.7\textwidth]{b2020_t1.png}
\end{figure}
\begin{footnotesize}
  Source: Individual responses from World Value  Survey, regressing compliance on individual characteristics, country and wave dummies. Income is displayed from lowest to highest
\end{footnotesize}
\end{frame}

\begin{frame}
\frametitle{Model setup} 
Basic structure of the model
\begin{itemize}\setlength\itemsep{3pt}
\item Share $e$ of elites who decides public goods and transfers, and non-elites
\item Among non-elites, $\mu_s$ of civic-minded citizens and $1-\mu_s$ of materialists
\begin{itemize}\setlength\itemsep{3pt}
\item All subject to taxes, with possible noncompliance ($n$), and receive utility from public goods 
\item Civic-minded: Gains if tax used on public goods ($G$), loss if used on transfers to elites ($B$)
\item Noncompliance: Incurs cost $C(n)$, which rises with fiscal capacity $\tau$ (easily detected)
\item Individual Utility: $\lambda>0$ for civic-minded citizens, $\lambda=0$ for materialists\vspace{-5pt}
\[\alpha G + b + w(1-(1-n)[t-\lambda\{G-eB\}-\tau C(n)])\] 
\end{itemize}\vspace{-10pt}
\item Revenues must be spent on public goods $G$, transfers to elites ($B$) and non-elites ($b$)
\begin{itemize}\setlength\itemsep{3pt}
\item For every $B$, share $0\leq\sigma\leq1$ must be spent on $b$. Thus $b=\sigma B$  (Higher $\sigma$ $\to$ cohesiveness)
\item Budget: $eB+(1-e)\sigma B = T-G $
\end{itemize}
\item Fiscal capacity is determined by physical investments but also `tax morale'
\begin{itemize}\setlength\itemsep{3pt}
  \item For materialists, only physical investments matter
\item For civic-minded person, public good$\uparrow \to$ tax morale $\uparrow$ (transfer$\uparrow\to$ tax morale $\downarrow$)  
\end{itemize}
\end{itemize}
\end{frame}

\begin{frame}
\frametitle{Optimal policy choice}  
Optimal policy choice
\begin{itemize}\setlength\itemsep{3pt}
\item Elites choose $t$, $G$, and $B$ to maximize the sum of individual utilities
\item Subject to cohesiveness of institutions ($\sigma$), fiscal capacity ($\tau$) and civic culture ($\mu$)
\item Whether revenues are spent on $G$ or $B$ depends on cohesiveness and civic culture
\begin{itemize}
\item Public goods are more likely to be provided in scenario where civic culture ($\mu\Uparrow$) and cohesiveness ($\sigma\Uparrow$) are dominant (and also raise fiscal capacity)
\item This comes from the fact that civic-minded persons are more likely to comply if revenues are used for public goods (and vice versa)
\item Furthermore, high cohesiveness makes transfers to elites expensive $\to$ make public goods relatively cheaper
\end{itemize}
\item Over time, additional public goods due to $(\mu, \sigma) \Uparrow$ $\Rightarrow$ benefits of being civic-minded $\Uparrow$
\begin{itemize}
\item Conversely, narrow range of public goods provision from $(\mu, \sigma) \Downarrow$ $\Rightarrow$ civic-mindedness $\Downarrow$
\end{itemize}
\item Initial state of civicness, cohesiveness + strength of common interest ($\alpha$) matters
\end{itemize}\medskip
\end{frame}

\begin{frame}
\frametitle{Core implications of the model} 
Cohesiveness of institutions and civic culture is complementary
\begin{itemize}\setlength\itemsep{3pt}
\item Cohesiveness of institutions are strengthened when more civic-minded persons induce high public goods over transfers to elites
\item People are more likely to be civic minded if institutional setups make transfers to elites costly (through cohesiveness)
\end{itemize}\medskip
These determine how common interests develop
\begin{itemize}\setlength\itemsep{3pt}
\item In a less cohesive, weakly civic-minded setting, external threat may not significantly increase common interest
\end{itemize}\medskip
Civic-mindedness and cohesivess can be fostered in more representative political setup
\begin{itemize}\setlength\itemsep{3pt}
\item Elites that more closely represents non-elites more likely to sustain power
\item Further increases returns for being civic-minded
\end{itemize}\medskip
\end{frame}

\begin{frame}
\frametitle{Potential extensions} 
High dependence on outside sources of funds
\begin{itemize}\setlength\itemsep{3pt}
\item Outside sources of funding reduces need for own-source tax revenues
\item This reduces importance of tax compliance and the role of civic culture
\item Sudden, exogenous changes to outside funds $\to$ Challenges in raising fiscal capacity 
\end{itemize}\medskip
Can elites determine civic-mindedness ($\mu$)?
\begin{itemize}\setlength\itemsep{3pt}
\item We assumed that $\mu$ is exogenously determined
\item But these can be potentially determined by elites
\item Room for extension: Making tax avoidance costly, Benefits of public services over time
\end{itemize}\medskip
\end{frame}

\begin{frame}
\frametitle{Takeaways: How these topics can be empirically studied} 
How can fiscal capacities be built 
\begin{itemize}\setlength\itemsep{3pt}
\item Evolution towards broad-based system (Gordon, Li 2009; Jensen 2022)
\item Efforts to gather more accurate information in developing countries (Balan et al 2021)
\end{itemize}\medskip
What increases tax morale and compliance?
\begin{itemize}\setlength\itemsep{3pt}
\item Increase benefits of compliance (tax breaks for formal firms) - (Benhassine 2018)
\item Tax perceptions and electoral outcomes (Christensen, Garfias 2021)
\end{itemize}\medskip
Other factors that contribute to state capacities
\begin{itemize}\setlength\itemsep{3pt}
\item Role of bureaucrats: Well-motivated, well-trained bureaucrats matter for effective public goods provision
\item Active studies on how to recruit and monitor them (Dal Bó et al 2013)
\end{itemize}
\end{frame}


\subsection{Rogowski, Gerring, Maguire, Cojocaru (2022): Public Infrastructure and Economic Development: Evidence from Postal Systems [student presentation]}
\begin{transitionframe}
\begin{center}
  \begin{huge}
    \textcolor{NTHU1}{%
Rogowski, Gerring, Maguire, Cojocaru (2022):\\ Public Infrastructure and Economic Development: Evidence from Postal Systems
    }
  \end{huge}
\end{center}
\end{transitionframe}




\section{Symptoms of Weak States}
\subsection{Shleifer, Vishny (1993): Corruption}
\begin{transitionframe}
\begin{center}
  \begin{huge}
    \textcolor{NTHU1}{%
  Shleifer, Vishny (1993):\\ Corruption
    }
  \end{huge}
\end{center}
\end{transitionframe}

\begin{frame}
\frametitle{Q: What conditions increase possibility of corruption}
Definition and setup on corruption
\begin{itemize}\setlength\itemsep{3pt}
\item Corruption: Sale by government officials of government property for personal gain
\begin{itemize}\setlength\itemsep{3pt}
\item Collecting bribes for providing permits, custom passage, bar competitor entry
\item Defense officials selling contracts for personal gain
\end{itemize}
\item In some cases, bribes \textit{may} account for a large fraction of national income
\begin{itemize}\setlength\itemsep{3pt}
\item In reality, really difficult to track. Most estimates may be underestimated
\end{itemize}
\end{itemize}\medskip
This paper provides framework on the organization and effects of corruption
\begin{itemize}\setlength\itemsep{3pt}
\item Compare different types of bribes (theft by officials vs not)  
\item Explain the distortionary effects of bribes 
\begin{itemize}\setlength\itemsep{3pt}
\item Bribe ``greases the oil?": Overcome cumbersome regulation etc
\item Slows down development, imposes additional costs
\end{itemize}
\end{itemize}
\end{frame}

\begin{frame}
\frametitle{Baseline model} 
Model setup
\begin{itemize}\setlength\itemsep{3pt}
\item Government good is up for sale (say, a passport or a license)
\item Official in charge can deny the sale of this good without risk of detection by its boss
\item Consider the official as monopolist: Maximize values of bribe for selling the good
\end{itemize}\medskip
Cost-benefit problem of the official
\begin{itemize}\setlength\itemsep{3pt}
\item Public good has a demand function $D(p)$, where $p$ is the official government price
\item For the official, no cost in \textit{producing} the good (already made by government)
\item Marginal cost of selling the good is different case by case
\begin{itemize}\setlength\itemsep{3pt}
\item Official turn over $p$ to government per sale (no theft): $p$ per sale $\to$ Buyer price $>p$ 
\item Officials turn nothing to government (theft): 0 marginal cost $\to$ Buyer price $<p$
\end{itemize}
\item Create a shortage of goods (first case) \& Reduce observance of law (second case)
\end{itemize}\medskip
\end{frame}

\begin{frame}
  \frametitle{Buyer price when corruption is present}
\begin{figure}
  \begin{subfigure}{0.49\textwidth}
  \includegraphics[keepaspectratio,width=\textwidth]{sv1993_f1_2.png}
  \caption{Without theft}
  \end{subfigure}
  \begin{subfigure}{0.49\textwidth}
  \includegraphics[keepaspectratio,width=\textwidth]{sv1993_f1_1.png}
  \caption{With theft}
  \end{subfigure}
\end{figure}
\end{frame}

\begin{frame}
\frametitle{What if you need multiple bribes to get a good?} 
Suppose you need more than one licenses 
\begin{itemize}\setlength\itemsep{3pt}
\item Two goods ($x_1, x_2$) are involved, with bribe and official prices of ($p_1, p_2$) and ($MC_1, MC_2$)
\item Officials may act independently or collude
\begin{itemize}
\item In former, each sets their own bribe $\to$ bribe prices are higher $\to$ underprovision of goods
\item In latter, keep bribe on one good down and profit off from bribe on the other
\end{itemize}
\item If a given good is supplied by different competing officials, level of bribes $\Downarrow$
\begin{itemize}
\item If someone asks for a bribe, go to the other official that do not ask for observance
\end{itemize}
\item The first two cases lead to underprovision and higher bribes compared to the third
\end{itemize}\medskip
What determines which one of the above structure are realized?
\begin{itemize}\setlength\itemsep{3pt}
\item How responsive is public pressure against corruption? (e.g active monitoring)
\item How active is the political competition (which leads to pressure against corruption)
\end{itemize}
\end{frame}

\begin{frame}
\frametitle{How does corruption further distort outcomes?} 
 Corruption is usually done in secret
\begin{itemize}\setlength\itemsep{3pt}
\item Thus, more actions need to be carried out by officials to avoid detection
\item May involve inducing substitutions into goods where bribes are hard to detection
\begin{itemize}
\item Results in higher bribes (to cover costs) and further underprovision of certain goods
\end{itemize}
\item These actions increase social costs,  much more than taxes
\end{itemize}\medskip
Other costs incurred due to secrecy
\begin{itemize}\setlength\itemsep{3pt}
\item Reluctance to change and innovation (keep outsiders off)
\item Reluctance to being open (hinder efforts for transparency)
\end{itemize}
\end{frame}

\begin{frame}
\frametitle{Updated analysis on corruption so far}
How to curb incentives for corruption among bureaucrats?
\begin{itemize}\setlength\itemsep{3pt}
\item Using data from public procurement (Best et al 2023)
\begin{itemize}
\item Breaks down prices charged in procurements to corruption and non-corrupt factors
\end{itemize}
\item Using bureaucrat placement (Khan et al 2019; Xu et al 2021)
\begin{itemize}\setlength\itemsep{3pt}
\item Finds that those assigned to home states are less effective
\item Other evidence that performance-based postings lead to better outcomes
\end{itemize}
\end{itemize}\medskip
How to increase monitoring against corruption? 
\begin{itemize}\setlength\itemsep{3pt}
\item Effects of media reports on capture of public funds (Reinikka, Svensson 2005)
\begin{itemize}
\item Provide information on local officials’ usage of funds $\to$ reduce wasted expenditures
\end{itemize}
\item Effects of reports on corrupt officials (Ferraz, Finnan 2008; Larreguy et al 2020)
\begin{itemize}
\item Publicize audit on local officials' performance $\to$ Accountability $\Uparrow$ ($\because$ electoral outcomes)
\end{itemize}
\end{itemize}
\end{frame}

\subsection{Dal Bó, Dal Bó, Di Tella (2006): Plata o Plomo? Bribe and Punishments in a Theory of Political Influence }
\begin{transitionframe}
\begin{center}
  \begin{huge}
    \textcolor{NTHU1}{%
Dal Bó, Dal Bó, Di Tella (2006):\\ Plata o Plomo? Bribe and Punishments in a Theory of Political Influence
    }
  \end{huge}
\end{center}
\end{transitionframe}

\begin{frame}
\frametitle{Q: How do state capture affect quality of public officials} 
In weak states, politicians are subject to punishments by private entities and bribes
\begin{itemize}\setlength\itemsep{3pt}
  \item Many private entities seek to influence politicians through bribes and coercion
\item Punishments and bribes complement: Threat of punishment makes bribing easier
\item Yet, there are things where models with only corruption cannot explain
\begin{itemize}
\item Why give officials more discretion and immunity?
\item When is it justified for the state to monopolize violence?
\end{itemize}
\end{itemize}\medskip
This paper updates framework of political influence to include bribes and coercion
\begin{itemize}\setlength\itemsep{3pt}
\item Investigates how possibility of coercion affects quality of public workers 
\item Explains why political immunities exist in contexts where private coercion is prevalent
\item Explains how discretion actually makes public work less attractive
\begin{itemize}\setlength\itemsep{3pt}
\item In model with bribes only, immunities are rarely justified (weaken monitoring mechanisms)
\item Also in the old model, discretion increases corruption and appeal of public work
\end{itemize}
\end{itemize}
\end{frame}

\begin{frame}
\frametitle{Model setup: Decision to join public work}
First stage: Citizens decide to enter public office or not 
\begin{itemize}\setlength\itemsep{3pt}

\item Citizens endowed with ability $a\in[0,\infty)$ (a probability distribution)
\item Alternative is to choose private sector (which gives wage $a$, equivalent to their ability)
\item Return from public work: public sector wage ($w$) + behavior of pressure group 
\item Citizens with $a\leq w$ apply (and $w$ types selected, assuming recruiters observe types)
\end{itemize} \medskip
Second stage: Pressure group has an opportunity to influence officials
\begin{itemize}\setlength\itemsep{3pt}
\item Public workers provide public good, whose quality positively correlates with ability
\item Have discretion to allocate ($\pi$) towards pressure group, with cost $c$ if caught
\item With probability $\gamma$, officials decide to accept bribes or face punishments
\begin{itemize}
  \item $1-\gamma$: Only punishment (no possibility to accept bribes) - when bribes are infeasible
\end{itemize}
\item Pressure groups can either bribe ($b$) or punish ($p$) at costs $\beta\Phi(b)$ and $\rho\Psi(p)$
\end{itemize}
\end{frame}

\begin{frame}
\frametitle{Analysis of payoffs} 
Decision on accepting (and handing out) bribes
\begin{itemize}\setlength\itemsep{3pt}
\item Officials facing pressure groups accept bribes (share $\gamma$) if \vspace{-5pt}
\[
w+b-c \geq w-p
\]
\vspace{-20pt}
\item Pressure groups sets $b$ and $p$ to maximize the following profits\vspace{-5pt}
\[
\Pi(b,p)=\gamma(\pi-\beta\Phi(b))-(1-\gamma)\rho\Psi(p) \ \ \text{\textit{subject to }}b\geq c-p
\]
\vspace{-20pt}
\item $\tilde{\pi}$: Minimal discretion required for pressure groups to influence officials (break-even)
\begin{itemize}
\item Inversely related to state capture: $\tilde{\pi}\Uparrow$ $\to$ Cost of bribing an official $\Uparrow$ $\to$ State capture  $\Downarrow$
\end{itemize}
\end{itemize}\medskip
If only bribes are available ($p=0$ for all cases)
\begin{itemize}\setlength\itemsep{3pt}
\item Solve $\gamma(\pi-\beta\Psi(b))$ subject to $b\geq c$ $\to$Bribe just enough to cover $c$, ($b=c$)
\item Official payoffs/skill is $w$, and $\tilde{\pi}^0=\beta\Phi(c)$
\item Drop in bribe costs ($\beta\downarrow$) $\to$ Higher state capture ($\tilde{\pi}\Downarrow$), no effect on quality of officials
\end{itemize}
\end{frame}

\begin{frame}
\frametitle{What happens when both bribes and coercion are possible?} 
Problem when both $b$ and $p$ are active
\begin{itemize}\setlength\itemsep{3pt}
\item Pay bribes just enough to cover $b=c-p$ (Bribes become cheaper)
\item The shortfall is made up with nonzero punishment $p^*$, such that $b+p^*=c$
\item $\tilde{\pi}^p$ satisfies $\frac{1-\gamma}{\gamma}\rho\Psi(p^*)+\beta\Phi(c-p^*)$ $\to$ Any discretion $\pi$ larger than this, then influence 
\item Resulting equilibrium payoff of official if pressure groups are active
\begin{itemize}\setlength\itemsep{3pt}
\item If official refuses bribe, payoff is $w-p^*$ (lower due to nonzero punishments)
\item If bribes are accepted, payoff is $w+b-c = w+c-p^* -c = w-p^*$
\end{itemize}
\item When pressure groups are inactive: payoff is $w>w-p^*$
\item Thus, the quality of bureaucrats fall in this scenario
\item Furthermore, $\tilde{\pi}^p<\tilde{\pi}^0$ $\to$ When threats are usable, more state capture!  
\end{itemize}
\end{frame}

\begin{frame}
\frametitle{Negative effects amplified with cheaper ways to influence officials} 
Exogenous changes in cost of influence (Suppose that $\beta$ and $\rho$ drops)
\begin{itemize}\setlength\itemsep{3pt}
\item Recall that $\tilde{\pi}^p=\frac{1-\gamma}{\gamma}\rho\Psi(p^*)+\beta\Phi(c-p^*)$ and $\Pi(b,p)=\gamma(\pi-\beta\Phi(b))+(1-\gamma)\rho\Psi(p)$
\item Fall in any one of $\beta$  and $\rho$ raises profits and decreases $\tilde{\pi}^p$
\begin{itemize}\setlength\itemsep{3pt}
\item The latter implies state capture will be more rampant 
\item Existing pressure groups ramp up activity + new pressure groups enter
\end{itemize}
\item If primarily driven by $\rho\Downarrow$, then payoff of officials (and quality) definitely falls
 \begin{itemize}\setlength\itemsep{3pt}
\item Usage of coercion increases in any case
\end{itemize}
\item If primarily driven by $\beta\Downarrow$, change to pay off ambiguous (but definitely not increase)
 \begin{itemize}\setlength\itemsep{3pt}
\item Existing groups shift away from coercion $\to$ Minimal effect on wages
\item New entrants that use both coercion and bribes $\to$ lower payoffs

\end{itemize}
\end{itemize}
\end{frame}

\begin{frame}
\frametitle{Implication on discretions}
How do discretion affect quality of politicians 
\begin{itemize}\setlength\itemsep{3pt}
\item If discretion $\pi<\tilde{\pi}^p$, less incentives for pressure groups to influence officials
\item Higher discretion could lead to more waste (more corrupt decision)
\item As more groups with coercion decrease payoffs, quality of officials worsened
\end{itemize}\medskip
Immunity of officials (in principle)
\begin{itemize}\setlength\itemsep{3pt}
\item Can insulate honest officials from outside pressures, mitigate false accusations
\item But can mask corrupt officials
\item In setting with rampant coercion from pressure groups
\begin{itemize}\setlength\itemsep{3pt}
\item There are greater effect of protection to officials with immunity
\item This can overweigh bad effects of masking corrupt officials
\end{itemize}
\item When legal system is effective, masking effects outweigh pressure effects
\end{itemize}
\end{frame}

\begin{frame}
\frametitle{Takeaways and discussions} 
Threat of coercion imposes negative effects on governance
\begin{itemize}\setlength\itemsep{3pt}
\item Reduces payoff for officials $\to$ lower quality of officials
\item More influence from private pressure groups $\to$ Increased state capture
\item Augments effects of bribes: Cheaper to bribe officials
\end{itemize}\medskip
Remaining questions
\begin{itemize}\setlength\itemsep{3pt}
\item Ultimately, need to improve legal capacity 
\begin{itemize}\setlength\itemsep{3pt}
\item To prevent pressure groups and detect corruption
\item But transitioning to better legal system is tough
\end{itemize}
\item Incentives of bureaucrats: How to attract incorruptible moral officials?
\begin{itemize}\setlength\itemsep{3pt}
\item Under what conditions can governments bring them in (what sort of incentives?)
\item Topic of Dal Bó et al (2013) in two weeks
\end{itemize}
\end{itemize}\medskip
\end{frame}


\subsection{Sanchez de la Sierra (2020): On the Origins of the State: Bandits and Taxation in Eastern Congo [Student Presentation]}
\begin{transitionframe}
\begin{center}
  \begin{huge}
    \textcolor{NTHU1}{%
Sanchez de la Sierra (2020):\\ On the Origins of the State: Bandits and Taxation in Eastern Congo [Student Presentation]
    }
  \end{huge}
\end{center}
\end{transitionframe}

\begin{frame}
\frametitle{Takeaways and remaining questions} 
Capable states are the foundation for economic and institutional development
\begin{itemize}\setlength\itemsep{3pt}
\item Able to collect taxes (or any resources) for common interests
\item Provision of essential public services
\item Protect market activities with legal measures
\item Role played by how much significance of common interest and civic movements 
\begin{itemize}
  \item Internal conflicts debilitate state capacity, external conflicts likely raise it
\end{itemize}
\item Weak states experience corruption and state capture (leading to further corruption)
\end{itemize}\medskip
How are we doing this empirically
\begin{itemize}\setlength\itemsep{3pt}
\item How do developing countries create policy instruments to enhance state capacity?
 \begin{itemize}
  \item Field and natural experiments on fiscal capacity
\end{itemize}
\item Studies on the bureaucrats implementing state operations
 \begin{itemize}
  \item How to motivate and monitor them?
\end{itemize}
\end{itemize}\medskip
\end{frame}


%%%%%%%%%%%
\end{document}


